\subsection{Problem 3: Phase Portraits}
In general, a phase portrait displays the evolution of the state of a dynamical system by plotting subsequent states in the same diagram.
For one-dimensional continuous systems the axes would usually be $\dot{x}$ and $x$.
For more than one dimension or state variable one has to choose them to be well representable in two dimensions.
In the system examined here we have an iterated map with two dimensions which are also the axes in the phase portrait plot.

If an attractive fixed point exists in the system, later iterations would be ever closer to that fixed point than previous iterations.
A stable orbit of length $n$ could not be distinguished from $n$ fixed points being present and would be needed to be further examined by the other methods employed in this report.
Chaotic conditions can be inferred when there is no discernible pattern in the phase portrait.
Pattern could here refer to either convergence or regularity.
